\section{Experimental Results}
\label{sec:results}

We validate all five frameworks on \mmlu capability data.
All experiments use 12 models with performance profiles across 57 subjects, constructed from published benchmark results.

\subsection{Framework 1: Topological Drift Detection}

\para{Persistence computation.}
We compute the Vietoris-Rips persistent homology of the full 12-model point cloud in $\sR^{57}$.
The persistence diagram reveals 12 $H_0$ features (connected components) with varying lifetimes and 0 $H_1$ features (loops), consistent with the small sample size relative to ambient dimension.

\para{Temporal drift.}
We partition models into early (2023: \gpt, \claude Opus, \llama 8B, PaLM-2, Mistral 7B, \mixtral) and late (2024: GPT-4o, Claude-3.5 Sonnet, \llama 70B, Qwen-2 72B, \gemini 1.5 Pro, Phi-3 Medium) cohorts.
\tabref{tab:drift} reports the distances between persistence diagrams of the two cohorts.

\begin{table}[t]
    \centering
    \caption{Topological drift between early (2023) and late (2024) model cohorts. The bottleneck distance provides a formal lower bound on the Hausdorff distance between the capability point clouds via the stability theorem.}
    \begin{tabular}{@{}lcc@{}}
        \toprule
        Metric & Value & Interpretation \\
        \midrule
        Wasserstein-1 distance & \textbf{1.77} & Total structural change \\
        Bottleneck distance & 0.68 & Maximum single-feature change \\
        Hausdorff lower bound & $\geq 0.34$ & Minimum geometric displacement \\
        $H_0$ features (full cloud) & 12 & Connected components \\
        $H_1$ features (full cloud) & 0 & No loops (small $n$) \\
        \bottomrule
    \end{tabular}
    \label{tab:drift}
\end{table}

The Wasserstein distance of 1.77 confirms genuine structural change between model generations.
By Theorem~\ref{thm:stability}, the Hausdorff distance between the point clouds is at least 0.34, establishing a formal lower bound on geometric displacement that no existing metric provides.

\subsection{Framework 2: Sheaf-Theoretic Composition}

\para{Graph structure.}
Our evaluation sheaf is defined on a graph with 63 nodes (1 global, 5 categories, 57 subjects) and 65 edges.
Categories comprise: STEM (17 subjects), Humanities (11), Social Sciences (11), Professional (10), Other (11).

\para{Consistency analysis.}
\tabref{tab:sheaf} reports consistency ratios for selected models.
Higher ratios indicate more heterogeneous profiles where category-level averaging loses substantial information.

\begin{table}[t]
    \centering
    \caption{Sheaf consistency ratios for selected models. Higher values indicate that standard averaging loses more information about the model's true capability distribution. Models with heterogeneous profiles (ratio $> 1$) cannot be faithfully evaluated by aggregated scores.}
    \begin{tabular}{@{}lcc@{}}
        \toprule
        Model & Consistency Ratio $\rho$ & Profile Type \\
        \midrule
        \gemini 1.5 Pro & \textbf{1.83} & Highly heterogeneous \\
        \mixtral & 1.72 & Heterogeneous \\
        GPT-4o & 1.15 & Moderately heterogeneous \\
        \gpt & 0.91 & Moderate \\
        \claude Opus & 0.76 & Relatively uniform \\
        \llama 8B & 0.48 & Uniform \\
        \bottomrule
    \end{tabular}
    \label{tab:sheaf}
\end{table}

\para{Key finding.}
\gemini 1.5 Pro (ratio 1.83) and \mixtral (ratio 1.72) exhibit highly heterogeneous profiles where category-level averaging discards substantial information.
For these models, reporting aggregate accuracy is formally misleading---the sheaf cohomology is non-trivial, indicating inconsistency between local and global evaluations.
Conversely, \llama 8B (ratio 0.48) has a sufficiently uniform profile that aggregation is faithful.

\subsection{Framework 3: Information-Geometric Manifold}

\para{Fisher-Rao distances.}
We compute pairwise \fr distances between all 12 models on the normalized capability simplex.
\tabref{tab:fisher} reports summary statistics and notable findings.

\begin{table}[t]
    \centering
    \caption{\fr distance statistics across all model pairs. The maximum distance (2.82) indicates substantial distributional difference, while geometrically close models with different raw accuracies reveal structure invisible to standard metrics.}
    \begin{tabular}{@{}lc@{}}
        \toprule
        Statistic & Value \\
        \midrule
        Maximum pairwise distance & \textbf{2.82} (Mistral 7B $\leftrightarrow$ GPT-4o) \\
        Mean pairwise distance & 1.46 \\
        Minimum pairwise distance & 0.31 (\claude Opus $\leftrightarrow$ \gpt) \\
        \bottomrule
    \end{tabular}
    \label{tab:fisher}
\end{table}

\para{Geometric structure.}
Multidimensional scaling (MDS) of the \fr distance matrix reveals clear separation between capability tiers with interesting intermediate structure.
Notably, \claude Opus and \gpt are geometrically close (distance 0.31) despite different raw accuracy values, indicating similar \emph{distributional structure} across subjects.
This finding is invisible to accuracy-based comparison and reveals shared architectural biases in capability allocation.

\para{Fisher information traces.}
Weaker models (Mistral 7B, \llama 8B) exhibit higher Fisher information traces, indicating more peaked capability distributions concentrated on fewer subjects.
These models are more sensitive to evaluation perturbation---a formal characterization of evaluation instability for weaker systems.

\subsection{Framework 4: Failure Mode Landscape}

\para{Phase transition.}
We compute persistent homology of the failure correlation graph at increasing thresholds.
\tabref{tab:failure} shows the progression of topological features.

\begin{table}[t]
    \centering
    \caption{Topological features of the failure landscape at increasing difficulty thresholds. A phase transition at $\tau = 0.8$ produces $H_1$ features (loops), indicating structurally distinct cyclic failure relationships that cannot be captured by hierarchical clustering.}
    \begin{tabular}{@{}ccccc@{}}
        \toprule
        Threshold $\tau$ & Failure Rate & $\beta_0$ (Components) & $\beta_1$ (Loops) & Interpretation \\
        \midrule
        0.5 & 3.9\% & 1 & 0 & Connected, no cycles \\
        0.6 & 9.8\% & 1 & 0 & Connected, no cycles \\
        0.7 & 21.1\% & 1 & 0 & Connected, no cycles \\
        0.8 & 54.8\% & 1 & \textbf{2} & Phase transition \\
        0.9 & 82.4\% & 1 & 3 & Rich cyclic structure \\
        \bottomrule
    \end{tabular}
    \label{tab:failure}
\end{table}

\para{Interpretation.}
The emergence of $H_1$ features at $\tau = 0.8$ represents a genuine phase transition: below this threshold, failure patterns form a tree-like hierarchy (subjects cluster by difficulty), but above it, cyclic failure relationships appear.
This indicates that at high difficulty thresholds, failure modes become structurally intertwined---models that fail on subject A tend to succeed on B, which tends to co-fail with C, which co-fails with A, forming non-trivial loops.

\para{Failure correlation structure.}
The failure correlation matrix exhibits clear block structure aligned with subject categories, confirming that failures are structured (non-random) and revealing inter-category failure dependencies invisible to per-subject analysis.

\subsection{Framework 5: Spectral Contamination Detection}

\para{Marchenko-Pastur analysis.}
For our performance matrix ($n=12$ models, $d=57$ subjects, $\gamma = d/n = 4.75$), the \marchpast bounds are $[\lambda_- = 0.29, \lambda_+ = 2.13]$.
\tabref{tab:spectral} summarizes the spectral analysis.

\begin{table}[t]
    \centering
    \caption{Spectral analysis results for contamination detection. Clean data produces 4 outlier eigenvalues reflecting genuine capability structure. Synthetic contamination produces additional outliers detectable via \tw statistics.}
    \begin{tabular}{@{}lcc@{}}
        \toprule
        Metric & Clean Data & Contaminated Data \\
        \midrule
        \marchpast bounds $[\lambda_-, \lambda_+]$ & \multicolumn{2}{c}{[0.29, 2.13]} \\
        Outlier eigenvalues ($> \lambda_+$) & 4 & 7 \\
        \tw statistic (largest $\lambda$) & \textbf{83.1} & 127.4 \\
        $p$-value & $< 10^{-10}$ & $< 10^{-10}$ \\
        Eigenvalue spacing & \multicolumn{2}{c}{Between Wigner and Poisson} \\
        \bottomrule
    \end{tabular}
    \label{tab:spectral}
\end{table}

\para{Clean data analysis.}
The clean performance matrix exhibits 4 eigenvalues exceeding $\lambda_+ = 2.13$, reflecting genuine capability structure (model tiers, category-specific strengths).
The \tw statistic of 83.1 is highly significant ($p < 10^{-10}$), confirming real structure beyond random noise.
The eigenvalue spacing lies between the Wigner surmise (correlated random matrix) and Poisson (independent entries), indicating moderate inter-task correlations.

\para{Contamination injection.}
We inject synthetic contamination by boosting 5 randomly selected subjects by 15\% for 4 models.
The contaminated matrix produces 3 additional outlier eigenvalues, and the corresponding eigenvectors' leverage scores correctly identify the contaminated subjects.
This demonstrates that spectral methods can detect contamination at a structural level, even when surface-level statistics remain plausible.

\subsection{Comparison to Standard Metrics}

\tabref{tab:comparison} summarizes the properties of our frameworks compared to standard evaluation approaches.

\begin{table}[t]
    \centering
    \caption{Comparison of mathematical properties across evaluation approaches. Our frameworks are the first to provide formal stability guarantees, composability analysis, and spectral contamination detection for LLM evaluation.}
    \resizebox{\textwidth}{!}{%
    \begin{tabular}{@{}lccccc@{}}
        \toprule
        Property & Accuracy & Elo Rating & \topdrift & \sheafeval & \speccontam \\
        \midrule
        Formal stability bound & \texttimes & \texttimes & \checkmark & --- & --- \\
        Composability guarantee & \texttimes & \texttimes & --- & \checkmark & --- \\
        Metric structure & \texttimes & partial & \checkmark & --- & \checkmark \\
        Contamination detection & \texttimes & \texttimes & --- & --- & \checkmark \\
        Multi-scale analysis & \texttimes & \texttimes & \checkmark & \checkmark & --- \\
        Distributional comparison & \texttimes & \texttimes & --- & --- & \checkmark \\
        \bottomrule
    \end{tabular}
    }
    \label{tab:comparison}
\end{table}
